\section{Introduction}
\label{sec:introduction}

\subsection{Motivation}
\label{sec:motivation}
\begin{itemize}
    \item Wide-area maritime surveillance: keeping a picture of vessel traffic over areas far larger than any coastal radar covers.
    \item Satellite-borne sensors (SAR, NRD, EO) cover such areas, but give position-only detections, with clutter, partial coverage per pass and long revisit times.
    \item AIS is cooperative and does not cover every vessel, which is why non-cooperative sensors and a tracker that combines their detections over time are needed. Here AIS is used as ground truth.
    \item Tracking (not just detecting) is needed to get identities, trajectories and a vessel count over time.
    \item This report is a specialization project that lays the groundwork for a Master's thesis in the same setting.
\end{itemize}

\subsection{Setting and scope}
\label{sec:scope}
\begin{itemize}
    \item One table, two columns: \emph{target setting} (the thesis) and \emph{this project}. Rows:
    \begin{itemize}
        \item Targets: point targets (vessels), position-only detections, no appearance features. Same in this project.
        \item Sensors: heterogeneous, asynchronous satellite sensors (SAR with clutter, NRD with per-measurement $R$ and almost no clutter, EO with constant $R$). This project: one simulated EO sensor, constant $R$ and $P_{\mathrm{D}}$.
        \item Revisit time: tens of minutes to hours. This project: 60--900~s.
        \item Area and density: large areas, many vessels, locally sparse. This project: up to 100~km radius, about 400 vessels.
        \item Coverage: partial per pass ($P_{\mathrm{D}}(x)=0$ outside the footprint). This project: static circular field of view, constant $P_{\mathrm{D}}$.
        \item Geometry: Earth curvature matters over long distances. This project: n-vector prediction and gating implemented.
        \item Motion: maneuvers over long intervals (e.g.\ ferries). This project: constant velocity only.
        \item Ground truth: AIS, interpolated to the scan times. Same in this project.
    \end{itemize}
    \item Out of scope, one sentence each: track-before-detect, extended objects, appearance/re-ID, end-to-end learned trackers, AIS as a sensor.
    \item Keep this subsection to about one page; Chapter~\ref{sec:mtt} starts from it.
\end{itemize}

\subsection{Problem statement and research questions}
\label{sec:rq}
\begin{itemize}
    \item Overall aim: a PMBM tracker that can be built on in the thesis for wide-area maritime tracking from satellite-type sensors.
    \item RQ1 (scale): Which reductions does a track-oriented PMBM need to run on realistic maritime scenes (hundreds of vessels, about 100~km radius, clutter up to $10^{-7}$ per m$^2$), and what do they cost in accuracy?
    \item RQ2 (geometry): Does exact spherical geometry (n-vector prediction and gating, \cite{b_gade_2024}) change tracking performance, and for which areas and revisit times should it be expected to?
    \item RQ3 (limits): Where, and why, does the resulting tracker fail as the revisit time grows towards satellite values?
\end{itemize}

\subsection{Contributions}
\label{sec:contributions}
\begin{itemize}
    \item A measurement simulator with real AIS tracks as ground truth, measurements and uniform clutter on the sphere.
    \item Integration of n-vector geometry into an existing PMBM implementation, including the transport of velocity and covariance along the great circle that \cite{b_gade_2024} leaves out.
    \item Clustering and recycling that make the PMBM run at about 400 vessels and realistic clutter.
    \item An evaluation of accuracy and run time against scale, clutter and revisit time, which identifies the limits that the thesis must address.
\end{itemize}

\subsection{Outline}
\label{sec:outline}
\begin{itemize}
    \item One sentence per chapter. Write last.
\end{itemize}
